\documentclass[11pt]{amsart}
\IfFileExists{lmodern.sty}%
  {\usepackage[T1]{fontenc}\usepackage{lmodern}\usepackage{microtype}}%
  {\usepackage[expansion=false]{microtype}}
\usepackage{amssymb,mathtools}
\usepackage{enumitem}
\emergencystretch=2.5em
\usepackage[colorlinks=true,linkcolor=blue,citecolor=blue,urlcolor=blue]{hyperref}
\hypersetup{pdftitle={A bound for the irrationality exponent of real elliptic periods},pdfauthor={Sean Litsey}}

\newtheorem{theorem}{Theorem}[section]
\newtheorem{lemma}[theorem]{Lemma}
\newtheorem{proposition}[theorem]{Proposition}
\newtheorem{corollary}[theorem]{Corollary}
\theoremstyle{definition}
\newtheorem{definition}[theorem]{Definition}
\theoremstyle{remark}
\newtheorem{remark}[theorem]{Remark}
\numberwithin{equation}{section}
\setcounter{tocdepth}{1}

\newcommand{\Z}{\mathbb Z}
\newcommand{\Q}{\mathbb Q}
\newcommand{\R}{\mathbb R}
\newcommand{\C}{\mathbb C}
\newcommand{\N}{\mathbb N}
\newcommand{\lcm}{\operatorname{lcm}}
\newcommand{\Ear}{E_{\mathrm{ar}}}
\newcommand{\Ean}{E_{\mathrm{an}}}
\newcommand{\cJ}{\mathcal J}
\newcommand{\cT}{\mathcal T}
\newcommand{\cF}{\mathcal F}
\newcommand{\cG}{\mathcal G}
\newcommand{\fa}{\mathcal I}
\newcommand{\tpsi}{\widetilde\psi}
\newcommand{\tg}{\widetilde g}
\newcommand{\tA}{\widetilde A}
\newcommand{\bb}{\bar b}
\newcommand{\CM}{\cite{ChudnovskyChudnovsky26}}

\begin{document}

\title{A bound for the irrationality exponent of real elliptic periods}

\author{Sean Litsey}
\email{sean.litsey.math@gmail.com}

\date{November 20, 2026}

\subjclass[2020]{Primary 11J82; Secondary 11J89, 11G05, 33E05}

\keywords{Irrationality exponent, elliptic periods, interpolation determinants, Weierstrass sigma function}

\begin{abstract}
We prove that every nonzero real period of an elliptic curve over $\Q$, with or without complex
multiplication, has irrationality exponent at most $3$. Equivalently, if $L\subset\C$ is a lattice whose
Weierstrass invariants $g_2(L)$ and $g_3(L)$ are rational, every nonzero real $\omega\in L$ satisfies
$\mu(\omega)\le3$. G.~and
D.~Chudnovsky proved $\mu(\omega)=2$ when the curve has complex multiplication. Their proof uses an
interpolation determinant of jets at centres on a two-dimensional grid of periods. Without complex
multiplication the only algebraic multiples of a real period inside the lattice are its integer multiples, so
the centres lie on a line. We show that the method still works on this one-dimensional grid. The arithmetic is
carried out over $\Z$ in the chart at the origin. A multiplier made periodic along the real period removes the
Gaussian factor in the transport of the determinant. The analytic estimate then costs a factor $K^2$ in place
of $K$, where $K$ is the number of centres. The two estimates still contradict each other for every exponent
larger than $3$. More generally, within the method an analytic cost $K^e$ gives exactly the exponents larger
than $e+1$. The bound is ineffective.
\end{abstract}

\maketitle
\tableofcontents

%% =====================================================================
\section{Introduction}\label{sec:intro}

\subsection{The result}
The \emph{irrationality exponent} $\mu(\xi)$ of a real number $\xi$ is the infimum of the real numbers $\nu$ such
that $|\xi-p/q|>q^{-\nu}$ for all integers $p$ and all sufficiently large integers $q$. It does not change when
$\xi$ is multiplied by a nonzero rational number. For irrational $\xi$ Dirichlet's theorem gives
$\mu(\xi)\ge2$.

\begin{theorem}\label{cor:curve}
Let $E$ be an elliptic curve over $\Q$, let $\eta$ be a nonzero invariant differential on $E$ defined over $\Q$,
and let $\omega$ be a nonzero real period of $\eta$. For every $\varepsilon>0$ there is an integer
$q_0=q_0(\varepsilon,\omega)$ such that
\[
 \Bigl|\omega-\frac pq\Bigr|>q^{-3-\varepsilon}
\]
for all integers $p$ and all integers $q\ge q_0$. In particular $\omega$ is irrational and
$2\le\mu(\omega)\le3$.
\end{theorem}

The curve may or may not have complex multiplication. For $E:y^2=x^3+Ax+B$ with $A,B\in\Q$ and $\eta=dx/(2y)$,
the periods are the points of a lattice in $\C$, which we now describe.

For a lattice $L\subset\C$ we write $g_2(L)=60\sum'\lambda^{-4}$ and $g_3(L)=140\sum'\lambda^{-6}$, the sums
running over $\lambda\in L\setminus\{0\}$. These are the invariants of the Weierstrass function $\wp_L$, which
satisfies $\wp_L'^2=4\wp_L^3-g_2(L)\wp_L-g_3(L)$. The curve $y^2=x^3+Ax+B$ with $4A^3+27B^2\ne0$ is uniformized
by the lattice $\Lambda$ with $g_2(\Lambda)=-4A$ and $g_3(\Lambda)=-4B$, through $x=\wp_\Lambda(z)$ and
$y=\wp_\Lambda'(z)/2$ \cite[Section~VI.5]{Silverman09}. The differential $dx/(2y)$ pulls back to $dz$, so its
periods form $\Lambda$. Conversely, every lattice $L$ with $g_2(L),g_3(L)\in\Q$ arises in this way from the curve
$y^2=x^3-\tfrac14g_2(L)x-\tfrac14g_3(L)$. Theorem~\ref{cor:curve} is therefore equivalent to the following
statement about lattices, which is the form we prove (Section~\ref{sec:proof}).

\begin{theorem}\label{thm:main}
Let $L\subset\C$ be a lattice with $g_2(L),g_3(L)\in\Q$, and let $\omega\in L$ be real and nonzero. For every
$\varepsilon>0$ there is an integer $q_0=q_0(\varepsilon,\omega)$ such that
$|\omega-p/q|>q^{-3-\varepsilon}$ for all integers $p$ and all integers $q\ge q_0$.
\end{theorem}

Theorem~\ref{cor:curve} applies to classical constants. The complete elliptic integral of the first kind is
\[
 K(k)=\int_0^{\pi/2}\frac{d\theta}{\sqrt{1-k^2\sin^2\theta}}=\frac{\pi}{2\,M\bigl(1,\sqrt{1-k^2}\bigr)}\qquad(0<k<1),
\]
where $M(a,b)$ is the arithmetic-geometric mean of Gauss.

\begin{corollary}\label{cor:K}
Let $0<k<1$ with $k^2\in\Q$. Then $\mu(K(k))\le3$, and equivalently $\mu\bigl(\pi/M(1,\sqrt{1-k^2})\bigr)\le3$.
\end{corollary}

\begin{proof}
Put $\lambda=k^2$. The curve $y^2=(1-x^2)(1-\lambda x^2)$ is smooth of genus one because $\lambda\ne0,1$, and it
has the rational point $(1,0)$. So it is an elliptic curve over $\Q$, and $dx/y$ is an invariant differential
defined over $\Q$. Its integral over the closed loop around the segment $[-1,1]$ equals
$2\int_{-1}^{1}dx/\sqrt{(1-x^2)(1-\lambda x^2)}=4K(k)$, a nonzero real period. Theorem~\ref{cor:curve} gives
$\mu(4K(k))\le3$, and $\mu$ does not change under multiplication by $4$.
\end{proof}

When the Legendre curve has complex multiplication, for instance $k^2=\tfrac12$ (the lemniscate case), the
exponent is $2$ by the work of G.~and D.~Chudnovsky \cite{ChudnovskyChudnovsky26}. Otherwise the corollary
gives, to our knowledge, the first explicit bound. For example, $k=\tfrac12$ gives the curve with
$j=2^8(\lambda^2-\lambda+1)^3/(\lambda^2(1-\lambda)^2)=35152/9$. The $j$-invariant of a curve with complex
multiplication is an algebraic integer, and $35152/9$ is not an integer. So this curve has no complex
multiplication, and $\mu(K(\tfrac12))\le3$. In the same way, Theorem~\ref{cor:curve} bounds by $3$ the
exponent of the real period $\Omega_E$ of every elliptic curve $E$ over $\Q$ in the tables of Cremona and the
LMFDB, for instance the curve \texttt{11a1} of smallest conductor.

The proof does not give $q_0(\varepsilon,\omega)$: the
bound is ineffective. A formalization in Lean~4 is intended to follow. The use of AI in this work is disclosed at
the end of the paper.

\subsection{Background}\label{ssec:background}
\emph{Rational approximation.} Roth proved that $\mu(\alpha)=2$ for every irrational algebraic number $\alpha$
\cite{Roth55}. For transcendental numbers the exact value of $\mu$ is known in few cases.

\emph{Elliptic periods.} Siegel proved that at least one of two basis periods of $\wp$ is transcendental when
$g_2$ and $g_3$ are algebraic \cite{Siegel32}. Under complex multiplication this makes every nonzero period
transcendental. Schneider proved that every nonzero period of such a $\wp$ is transcendental, with or without
complex multiplication \cite{Schneider37}. He also proved that if $\tau$ and the modular invariant $j(\tau)$
are both algebraic, then $\tau$ is imaginary quadratic \cite{Schneider37}. We use this second theorem in Remark~\ref{rem:noCM}. Both
results, and their later history, are described in Waldschmidt's survey \cite[Section~3.1]{Waldschmidt08}.

\emph{Measures.} Fel'dman gave measures of transcendence for the numbers in the theorems of Siegel and Schneider
\cite{Feldman51,Feldman74}. Masser determined the linear relations among $1$, the periods, the quasi-periods and
$2\pi i$ of one Weierstrass function, and he gave measures of linear independence for them \cite{Masser75}.
Philippon and Waldschmidt proved effective lower bounds for linear forms in logarithms on commutative algebraic
groups. They include the first such measures for elliptic logarithms that hold without complex multiplication
\cite{PhilipponWaldschmidt88,PhilipponWaldschmidt89}. Hirata-Kohno and David refined these bounds
\cite{HirataKohno90,HirataKohno91,David95}. As recorded in \cite[Section~4.3]{Waldschmidt08}, it follows from
work of David, Hirata-Kohno and Gaudron that these nonvanishing linear forms are not Liouville numbers. So
$\mu(\omega)$ is finite for every nonzero real period $\omega$, with an exponent that depends on the curve. We
have not found a numerical value of such an exponent in the literature. For curves without complex
multiplication, Theorem~\ref{cor:curve} gives, to our knowledge, the first explicit bound, and it is uniform over
all curves over $\Q$.

\emph{Recent results.} Since September 2026 several results of exponent $2$ have appeared, some obtained with AI
systems. In October 2026 OpenAI proved that $\mu(\pi)=2$, with a Lean formalization \cite{OpenAIPi2026}.
Liu, Jiang and Zhang proved $\mu(\log r)=2$ for rational $r>0$, $r\ne1$ \cite{LiuJiangZhang26}.
G.~V.~Chudnovsky proved $\mu(\varpi)=2$ for the lemniscate period
$\varpi=2\int_0^1(1-t^4)^{-1/2}\,dt$ \cite{Chudnovsky26}. G.~V.~and D.~V.~Chudnovsky extended this to every
nonzero real period of an elliptic curve over $\Q$ with complex multiplication, for an invariant differential
defined over $\Q$. They also proved $\mu(\varpi^n)=O(n\log 2n)$ \CM. The paper \CM{} is a preprint. Its results
that we use are stated exactly where they are used, in Theorem~\ref{thm:epa}, Lemma~\ref{lem:trunc} and
Lemma~\ref{lem:collision}.

\subsection{What the method needs without complex multiplication}\label{ssec:need}
The proof of \CM{} works on $E\times\mathbb G_a^m$. It takes $m$ rational approximations $p_i/q_i$ to a real
period $\omega$, and a set of $K$ periods $\gamma_j=\lambda_j\omega\in\Lambda$. The centres are the points
$(\lambda_jp_1/q_1,\dots,\lambda_jp_m/q_m)$ of the additive factor. The determinant $\Delta$ is formed from the
weighted jets of a space of polynomials at these centres. Three estimates are compared. An interpolation
theorem with separated weights shows that $\Delta\ne0$. Clearing denominators
gives a lower bound for $|\Delta|$, since $\lambda_j$ lies in an order of an imaginary quadratic field. An
analytic upper bound comes from the periodicity of the elliptic coordinates and a collision estimate. The
centres $\lambda_j$ run over a rectangle $\{a+bi\sqrt d\}$ in the order, so they lie in a disc of radius of order
$\sqrt K$. The multiplier that clears the poles grows like $\exp(c|z|^2)$, and the analytic estimate costs a
term of order $K/w_0$, where $w_0$ is the source weight of the elliptic coordinates. Apart from the
irrationality of $\omega$, complex multiplication enters \CM{} only through this two-dimensional grid.

\begin{remark}\label{rem:noCM}
Let $\Lambda$ have algebraic invariants and no complex multiplication, and let $\Omega$ generate
$\Lambda\cap\R$. If $\alpha$ is algebraic and $\alpha\Omega\in\Lambda$, then $\alpha\in\Z$. Indeed, write
$\Lambda=\Z\Omega+\Z\Omega'$ and $\alpha\Omega=a\Omega+b\Omega'$. If $b\ne0$, then $\tau=\Omega'/\Omega$ is
algebraic. Since $j(\tau)$ is algebraic, Schneider's theorem makes $\tau$ imaginary quadratic, and $\Lambda$ has
complex multiplication. So $b=0$ and $\alpha=a$.
\end{remark}

So without complex multiplication the multipliers $\lambda_j$ must be integers. We take the centres
$jp_i/q_i$ with $0\le j<K$. They lie on a line, at distance up to $K\Omega$ from the origin. The multiplier then costs a term of order $K^2/w_0$. The
question is whether the parameters still close, and at which exponent.

\subsection{New ideas}
We keep the construction of \CM{} and its nonvanishing theorem. The proof rests on the following ideas.
\begin{enumerate}[leftmargin=*]
\item \emph{A one-dimensional real grid.} The periods $j\Omega$, $0\le j<K$, replace the two-dimensional grid of
 complex multiplication, and the centres are $c_{ji}=jp_i/q_i$. These centres are distinct in each coordinate,
 which is all that the nonvanishing theorem of \CM{} asks (Theorem~\ref{thm:epa}). The entries lie in $\Q$, and
 the lower bound needs only $|n|\ge1$ for a nonzero integer $n$.
\item \emph{The origin chart for every curve.} For $y^2=x^3+Ax+B$ with $A,B\in\Z$ we use the chart of \CM{} at
 the origin, $u=1/(x-a)$, $t=y/(x-a)^2$, with any integer $a$ such that $a^3+Aa+B\ne0$. Its elliptic logarithm
 has the denominators $n$ only. The whole arithmetic part of \CM{} uses nothing else. We give it over $\Z$
 (Proposition~\ref{prop:arith}). No torsion point and no endomorphism enters.
\item \emph{The periodised multiplier.} Let $\psi(z)=\sigma(z-z_Q)\sigma(z+z_Q)$ be the sigma-type multiplier of
 \CM{} and let $\eta_\Omega$ be the quasi-period of $\Omega$. The function
 $\tpsi(z)=\psi(z)\exp(-\eta_\Omega z^2/\Omega)$ is periodic with period $\Omega$. In the transport of the
 determinant it removes the factor $\prod_j|\psi(j\Omega)|$, which is a Gaussian in $j$
 (Lemma~\ref{lem:transport}). Legendre's relation gives its growth across the strip,
 $\log|\tpsi(z)|\le c+\chi(|\operatorname{Im}z|+Y)^2$ for a constant $c$, with $\chi=2\pi/(\Omega Y)$ and $Y$ the
 height of the period parallelogram (Lemma~\ref{lem:growth}).
\item \emph{The analytic cost $K^2/w_0$ and the exponent $3$.} The collision estimate needs a disc around all
 centres, of radius of order $K$. On it the multiplier costs a term $C_\Lambda K^2/(2w_0)$, with an explicit
 constant $C_\Lambda$ of the lattice (Theorem~\ref{thm:analytic}). With explicit rational choices of all
 parameters the comparison closes for every exponent $\nu>3$ (Section~\ref{sec:params}).
\item \emph{The threshold $e+1$.} If the analytic cost is $K^e/w_0$, the parameters close exactly when
 $\nu>e+1$ (Proposition~\ref{prop:threshold}). The case $e=1$ is the exponent $2$ of complex multiplication and
 of $\pi$. The case $e=2$ is the exponent $3$ of this paper.
\end{enumerate}

\subsection{Constants, organization and notation}
The exponent $3$ is the only constant in Theorem~\ref{thm:main}. The constants of the proof, such as
$C_\Lambda$, depend on the lattice and do not enter the result.

Section~\ref{sec:chart} reduces Theorem~\ref{thm:main} to an integral model and constructs the chart and the
multiplier. Section~\ref{sec:det} defines the determinant and states the nonvanishing theorem.
Sections~\ref{sec:arith} and~\ref{sec:analytic} prove the lower and upper bounds. Section~\ref{sec:params}
chooses the parameters and Section~\ref{sec:proof} proves the theorem. Section~\ref{sec:three} explains why the
exponent is $3$ and what would give $2$.

Logarithms are natural, $\N$ contains $0$, and $\log^+x=\max(0,\log x)$. For a series $f$ in $T$ and
$V=(V_1,\dots,V_m)$, $[T^sV^\beta]f$ is the coefficient of $T^s\prod_iV_i^{\beta_i}$. For
$\alpha,\beta\in\N^m$ we write $\beta\le\alpha$ coordinatewise, $|\alpha|=\sum_i\alpha_i$ and
$\binom\alpha\beta=\prod_i\binom{\alpha_i}{\beta_i}$, which is $0$ unless $\beta\le\alpha$.

%% =====================================================================
\section{The lattice, the origin chart and the multiplier}\label{sec:chart}

\subsection{Reduction to an integral model}
The first lemma replaces $L$ by a lattice whose curve has integer coefficients.

\begin{lemma}\label{lem:integral}
Let $L$ be a lattice with $g_2(L),g_3(L)\in\Q$. There is a positive integer $d$ such that $\Lambda=d^{-1}L$
satisfies $g_2(\Lambda)=-4A$ and $g_3(\Lambda)=-4B$ with $A,B\in\Z$ and $4A^3+27B^2\ne0$. The functions
$x=\wp_\Lambda$ and $y=\wp'_\Lambda/2$ satisfy $y^2=x^3+Ax+B$ and $dx/(2y)=dz$.
\end{lemma}

\begin{proof}
For $c\in\C^\times$ one has $g_2(cL)=c^{-4}g_2(L)$ and $g_3(cL)=c^{-6}g_3(L)$. Let $s$ be a common denominator
of $g_2(L)$ and $g_3(L)$ and put $d=2s$. Then $d^4g_2(L)/4$ and $d^6g_3(L)/4$ are integers, and we define
$A=-d^4g_2(L)/4$ and $B=-d^6g_3(L)/4$. The discriminant $g_2^3-27g_3^2$ of a lattice is nonzero, and for
$\Lambda$ it equals $-16(4A^3+27B^2)$. The differential equation of $\wp_\Lambda$ gives
$y^2=\wp'^2_\Lambda/4=x^3+Ax+B$, and $dx/(2y)=\wp'_\Lambda\,dz/\wp'_\Lambda=dz$.
\end{proof}

A nonzero real $\omega\in L$ gives the nonzero real period $\omega/d\in\Lambda$. The group $\Lambda\cap\R$ is
discrete and nonzero, so $\Lambda\cap\R=\Z\Omega$ for a unique $\Omega>0$. Thus $\omega/d=n\Omega$ with a
nonzero integer $n$, and $\omega=N\Omega$ with $N=dn$. Integer multiples do not change the exponent, in the following form.

\begin{lemma}\label{lem:scale}
Let $N\ne0$ be an integer and $\nu'<\nu$. If infinitely many pairs $(p,q)$ with $q\ge1$ satisfy
$|N\xi-p/q|\le q^{-\nu}$, then infinitely many pairs $(p',q')$ with $q'\ge1$ satisfy
$|\xi-p'/q'|\le q'^{-\nu'}$.
\end{lemma}

\begin{proof}
For bounded $q$ only finitely many $p$ satisfy the inequality, so the denominators $q$ are unbounded. Put
$p'=p\operatorname{sgn}N$ and $q'=|N|q$. Then $|\xi-p'/q'|\le q^{-\nu}/|N|$. This is at most
$(|N|q)^{-\nu'}$ as soon as $q^{\nu-\nu'}\ge|N|^{\nu'-1}$, which holds for all large $q$.
\end{proof}

From now on we fix $A,B\in\Z$ with $4A^3+27B^2\ne0$, the cubic $f(X)=X^3+AX+B$, the lattice $\Lambda$ of
Lemma~\ref{lem:integral} and the generator $\Omega>0$ of $\Lambda\cap\R$. We choose $\Omega'\in\Lambda$ with
$\Lambda=\Z\Omega+\Z\Omega'$ and $Y=\operatorname{Im}\Omega'>0$.

\subsection{The origin chart}
This chart is the one of \CM. Fix an integer $a$ with $f(a)\ne0$. One of $0,1,2,3$ will do, since a cubic has at
most three roots. Put
\[
 u=\frac1{x-a},\qquad t=\frac y{(x-a)^2}.
\]
Multiplying $y^2=f(x)$ by $u^4$ gives
\begin{equation}\label{eq:quartic}
 t^2=F(u),\qquad F(U)=U+3aU^2+(3a^2+A)U^3+f(a)U^4\in\Z[U].
\end{equation}
The decisive facts are $F(0)=0$ and $F'(0)=1$.

\begin{proposition}[{\cite[Proposition~3.1]{ChudnovskyChudnovsky26}}]\label{prop:chart}
There is a unique $U(T)\in\Z[[T]]$ with $U(0)=0$ and $F(U(T))=T^2$. The series $-1/F'(U(T))$ lies in $\Z[[T]]$.
The series $L(T)\in\Q[[T]]$ with $L(0)=0$ and $L'(T)=-1/F'(U(T))$ satisfies $n[T^n]L\in\Z$ for $n\ge1$.
\end{proposition}

\begin{proof}
The linear coefficient of $F$ is $1$, so $F$ has a compositional inverse $G\in\Z[[S]]$ with $G(0)=0$. Put
$U(T)=G(T^2)$. Any solution $W$ with $W(0)=0$ satisfies $W=G(F(W))=G(T^2)$, which gives uniqueness. The
constant term of $F'(U(T))$ is $1$, so its inverse lies in $\Z[[T]]$. If $-1/F'(U(T))=\sum_kv_kT^k$, then
$[T^n]L=v_{n-1}/n$.
\end{proof}

Next we relate the chart to the uniformization. The function $u=1/(\wp_\Lambda-a)$ vanishes at the points of
$\Lambda$. Let $Q=(a,y_Q)$ be a point of the curve with $x(Q)=a$, and let $z_Q$ be an elliptic logarithm of $Q$.
Since $y_Q^2=f(a)\ne0$, the point $Q$ is not of order $2$, so the cosets $z_Q+\Lambda$ and $-z_Q+\Lambda$ are
distinct and $\wp'_\Lambda(z_Q)\ne0$. Hence $u$ has simple poles exactly at $\pm z_Q+\Lambda$. The function
$t=\tfrac12\wp'_\Lambda u^2$ has double poles there. At $z=0$ the triple pole of $\wp'_\Lambda$ meets the
zero of order four of $u^2$, so $t$ vanishes at the points of $\Lambda$ and has no other poles.

\begin{proposition}[{\cite[Proposition~3.3]{ChudnovskyChudnovsky26}}]\label{prop:germ}
Write $\tau(z)=t(z)$ near $z=0$. Then $\tau'(0)=-1$, and as convergent germs at $z=0$
\[
 u(z)=U(\tau(z)),\qquad L(\tau(z))=z.
\]
\end{proposition}

\begin{proof}
Near $0$ one has $\wp_\Lambda=z^{-2}+O(z^2)$ and $\wp'_\Lambda=-2z^{-3}+O(z)$. So $u=z^2+O(z^4)$ and
$t=-z+O(z^3)$. Thus $\tau$ is a local coordinate, and $u=W(\tau)$ for a convergent series $W$ with $W(0)=0$.
The identity $t^2=F(u)$ gives $F(W(T))=T^2$, and Proposition~\ref{prop:chart} gives $W=U$. From
$x=a+1/u$ and $y=t/u^2$ we get $dz=dx/(2y)=-du/(2t)$. Differentiating $t^2=F(u)$ gives $2t\,dt=F'(u)\,du$,
hence $dz=-dt/F'(u)=L'(\tau)\,d\tau$. Both $z$ and $L(\tau(z))$ vanish at $0$, so they are equal.
\end{proof}

\subsection{The periodised multiplier}
Let $\sigma$ be the Weierstrass sigma function of $\Lambda$. We use the classical
facts in the form of \cite[Section~23.2]{DLMF}. The function $\sigma$ is entire with simple zeros exactly on
$\Lambda$. For $\lambda\in\Lambda$ there are $\varepsilon_\lambda\in\{\pm1\}$ and $\eta_\lambda\in\C$ with
\begin{equation}\label{eq:quasi}
 \sigma(z+\lambda)=\varepsilon_\lambda e^{\eta_\lambda(z+\lambda/2)}\sigma(z),
\end{equation}
and $\lambda\mapsto\eta_\lambda$ is additive. Legendre's relation reads
\begin{equation}\label{eq:legendre}
 \eta_\Omega\Omega'-\eta_{\Omega'}\Omega=2\pi i.
\end{equation}
In the notation of \cite{DLMF}, $\Omega=2\omega_1$, $\Omega'=2\omega_3$, $\eta_\Omega=2\eta_1$ and
$\eta_{\Omega'}=2\eta_3$. Then \eqref{eq:quasi} is (23.2.17) and \eqref{eq:legendre} is four times (23.2.14).

Put
\[
 \psi(z)=\sigma(z-z_Q)\,\sigma(z+z_Q),\qquad \tpsi(z)=\psi(z)\exp\Bigl(-\frac{\eta_\Omega}{\Omega}z^2\Bigr).
\]
The function $\psi$ is the multiplier of \CM. The factor $\exp(-\eta_\Omega z^2/\Omega)$ is new.

\begin{lemma}\label{lem:mult}
\begin{enumerate}[label=(\alph*)]
\item $\psi(z+\lambda)=e^{\eta_\lambda(2z+\lambda)}\psi(z)$ for $\lambda\in\Lambda$, and
 $\tpsi(z+\Omega)=\tpsi(z)$.
\item $\tpsi(z+n\Omega')=\tpsi(z)\exp\bigl(-\tfrac{2\pi i}{\Omega}(2nz+n^2\Omega')\bigr)$ for $n\in\Z$.
\item The functions $\tpsi$, $\tpsi u$ and $\tpsi^2t$ are entire, and $\tpsi(0)=\psi(0)=-\sigma(z_Q)^2\ne0$.
\end{enumerate}
\end{lemma}

\begin{proof}
(a) Apply \eqref{eq:quasi} at $z-z_Q$ and $z+z_Q$. The signs give $\varepsilon_\lambda^2=1$ and the exponents add
to $\eta_\lambda(2z+\lambda)$. For $\lambda=\Omega$ the exponential factor of $\tpsi$ changes by
$\exp(-\eta_\Omega(2z+\Omega))$, which cancels.

(b) With $\lambda=n\Omega'$ the exponent becomes
$(n\eta_{\Omega'}-n\eta_\Omega\Omega'/\Omega)(2z+n\Omega')$. By \eqref{eq:legendre} the first factor is
$-2\pi in/\Omega$.

(c) The zeros of $\psi$ are simple and lie exactly on $\pm z_Q+\Lambda$. They cancel the simple poles of $u$ and
the double poles of $t$. Finally $z_Q\notin\Lambda$ and $\sigma$ is odd.
\end{proof}

For integers $D\ge1$ and $k,b\ge0$ with $k+2b\le D$ put
\[
 \tg_{D,k,b}=\tpsi^{D}u^kt^b=\tpsi^{D-k-2b}(\tpsi u)^k(\tpsi^2t)^b.
\]
It is entire and $\Omega$-periodic. Let $S_*\ge1$ be the maximum of $1$ and the suprema of $|\tpsi|$,
$|\tpsi u|$ and $|\tpsi^2t|$ on the rectangle $\{x+iy:0\le x\le\Omega,\ 0\le y\le Y\}$. Put
\[
 \chi=\frac{2\pi}{\Omega Y}.
\]

\begin{lemma}\label{lem:growth}
For all $z\in\C$ and all admissible $(D,k,b)$,
\[
 |\tg_{D,k,b}(z)|\le S_*^D\exp\bigl(D\chi(|\operatorname{Im}z|+Y)^2\bigr).
\]
\end{lemma}

\begin{proof}
Write $z=z_1+n\Omega'+l\Omega$ with $n,l\in\Z$, $0\le\operatorname{Re}z_1<\Omega$ and
$0\le y_1=\operatorname{Im}z_1<Y$. The functions $u,t$ are $\Lambda$-periodic and $\tpsi$ is
$\Omega$-periodic. By Lemma~\ref{lem:mult}(b), each of $\tpsi$ and $\tpsi u$ is multiplied by a factor of
modulus $\exp(\tfrac{2\pi}\Omega(2ny_1+n^2Y))$ in passing from $z_1$ to $z$, and $\tpsi^2t$ by its square. The
product defining $\tg_{D,k,b}$ therefore acquires this factor to the power $(D-k-2b)+k+2b=D$. At $z_1$ the
product has $D-b\le D$ factors, each bounded by $S_*$.

It remains to bound $2ny_1+n^2Y$ by $Y^{-1}(|y|+Y)^2$, where $y=\operatorname{Im}z=y_1+nY$. If $n\ge0$, then
$2ny_1+n^2Y\le(n+1)^2Y$ and $(n+1)Y\le y+Y$. If $n<0$, then $2ny_1+n^2Y\le n^2Y$, and
$|n|Y=y_1-y<Y+|y|$. In both cases the factor is at most $\exp(\chi(|y|+Y)^2)$.
\end{proof}

%% =====================================================================
\section{The interpolation determinant}\label{sec:det}

This section sets up the determinant of \CM{} for the one-dimensional grid, and states its nonvanishing.

\subsection{Weights, rows and columns}
Fix integers $m,K\ge1$, positive rationals $w_0,v_0$ and $\theta\in(0,1)$, and a cutoff $H>0$. Fix integers
$p_i\ne0$ and $q_i\ge2$ for $1\le i\le m$, and put $w_i=\lceil\log q_i\rceil$ and $W(\beta)=\sum_iw_i\beta_i$
for $\beta\in\N^m$. The target monomials and their number are
\[
 \cJ_H=\{(s,\beta)\in\N^{1+m}:v_0s+W(\beta)/\theta<H\},\qquad J_H=|\cJ_H|,\qquad M=KJ_H.
\]
The source space is
\[
 \cF_H=\operatorname{span}_\C\{u^kt^bZ^\alpha:\ k\in\N,\ b\in\{0,1\},\ \alpha\in\N^m,\
 2w_0(k+2b)+W(\alpha)\le H\}.
\]
The centres are
\[
 c_{ji}=j\,\frac{p_i}{q_i}\qquad(0\le j<K,\ 1\le i\le m).
\]
For each $i$ they are distinct as $j$ varies, since $p_i\ne0$. For $P\in\cF_H$, centre $j$ and
$(s,\beta)\in\cJ_H$, the coefficient
\begin{equation}\label{eq:packet}
 [T^sV^\beta]\,P\bigl(U(T),T,c_{j1}+V_1+L(T),\dots,c_{jm}+V_m+L(T)\bigr)
\end{equation}
is the $(s,\beta)$-jet of $P$ at the centre $j$, read in the chart of Proposition~\ref{prop:chart}. Choose $M$
monomials $P_c=u^{k_c}t^{b_c}Z^{\alpha_c}$ of $\cF_H$, so that
\begin{equation*}
 b_c\in\{0,1\},\qquad 2w_0(k_c+2b_c)+W(\alpha_c)\le H.\tag{S}\label{eq:S}
\end{equation*}
The matrix $A$ has rows $r=(j,s,\beta)$ with $0\le j<K$ and $(s,\beta)\in\cJ_H$, and columns $c$. Its entries
are the coefficients \eqref{eq:packet} for $P=P_c$. Taylor expansion in $V$ gives
\begin{equation}\label{eq:entry}
 A_{(j,s,\beta),c}=\binom{\alpha_c}{\beta}[T^s]\Bigl(U(T)^{k_c}T^{b_c}\prod_i\bigl(c_{ji}+L(T)\bigr)^{
 \alpha_{ci}-\beta_i}\Bigr).
\end{equation}
We write $\Delta=\det A$ and
\[
 \bb=\frac1{MH}\sum_{(j,s,\beta)}W(\beta).
\]
The target cutoff gives $W(\beta)<\theta H$ for every row, so $0\le\bb<\theta$.

\subsection{Nonvanishing}
The nonvanishing theorem of \CM{} is a separated-weight interpolation theorem in the chart at the origin. We
state it as it is proved there. For $n\in\N$ and $s>0$ put
\[
 A_n(s)=2(n+2)^{n+2}\Bigl(1+\frac{n+2}s\Bigr)^{n+2}.
\]

\begin{theorem}[{\cite[Corollary~13.5]{ChudnovskyChudnovsky26}}]\label{thm:epa}
Let $A,B,a$ be integers with $4A^3+27B^2\ne0$ and $f(a)\ne0$. Let $m,K\ge1$ and let $w_0,v_0,\theta$ be
positive rationals with
\begin{equation}\label{eq:epa-base}
 \theta<1,\qquad K\frac{w_0}{v_0}\theta^m<1,\qquad K\theta^m<1.
\end{equation}
Choose a rational $\sigma_0>0$ with
\[
 (1+3\sigma_0)^{m+1}K\frac{w_0}{v_0}\theta^m<1,\qquad(1+3\sigma_0)^mK\theta^m<1,\qquad(1+\sigma_0)\theta<1.
\]
Put $\varepsilon_0=\min\{\sigma_0,\sigma_0v_0/(4(1+\sigma_0)w_0)\}$,
$C_4=A_m(\sigma_0)+4\cdot2^{m+3}A_{m+1}(\varepsilon_0)$ and $q_*=\theta^m\min(1,w_0)/\max(1,v_0)$. Let
$w_1,\dots,w_m$ be positive rationals with
\begin{equation}\label{eq:epa-thresholds}
 w_i>\max\Bigl\{1,\frac{C_4}{q_*}\prod_{1\le l<i}w_l\Bigr\}\qquad(1\le i\le m).
\end{equation}
There is a positive integer $D_{\mathrm{EPA}}$, depending on these data but not on any centres, with the
following property. For every array $(c_{ji})\in\C^{K\times m}$ such that $j\mapsto c_{ji}$ is injective for
each $i$, there is $H_0$ such that for every integer $H\ge H_0$ divisible by $D_{\mathrm{EPA}}$ the map from
$\cF_H$ to $\prod_j\C^{\cJ_H}$ given by the coefficients \eqref{eq:packet} is surjective.
\end{theorem}

The theorem is stated in \CM{} for every curve $y^2=x^3+Ax+B$ with $A,B\in\Z$ and $4A^3+27B^2\ne0$, and for
every array whose coordinates are injective in $j$. Its proof, in \cite[Sections~4--13]{ChudnovskyChudnovsky26},
uses the elliptic curve through four facts only. The quartic $F$ is squarefree, which follows from
$4A^3+27B^2\ne0$. Its leading coefficient $f(a)$ is nonzero. The differential $dx/(2y)$ has no algebraic
primitive on any finite cover. Finally $F'(0)=1$. None of these uses an endomorphism. The centres enter through
the injectivity of each coordinate, and in one place through their distinctness as points of $\C^m$. Complex
multiplication is mentioned in those sections only to fix the curve. So Theorem~\ref{thm:epa} applies to every
curve and to the collinear centres $jp_i/q_i$.

\begin{remark}\label{rem:line}
The collinear centres do impose a condition, which the thresholds already contain. They lie on the line
$u=t=0$, $Z_i=\lambda p_i/q_i$. Its pole divisor has degree $1/w_1$ in the weighted sense of \CM, and it meets
each of the $K$ centres with contact $\theta/w_m$. The curve inequality behind Theorem~\ref{thm:epa} therefore
contains $(1+\sigma_0)K\theta\le w_m/w_1$. For $m\ge2$ this follows from \eqref{eq:epa-thresholds}, since
$w_m/w_1>C_4/q_*\ge\theta^{-m}>K>(1+\sigma_0)K\theta$. For $m=1$ it follows from the condition
$(1+3\sigma_0)K\theta<1$ on $\sigma_0$.
\end{remark}

If the coefficient map is surjective, the matrix with all $M$ rows and with the columns indexed by all
monomials of $\cF_H$ has rank $M$. Some choice of $M$ columns then gives $\Delta\ne0$. These columns satisfy
\eqref{eq:S}.

%% =====================================================================
\section{The arithmetic lower bound}\label{sec:arith}

The entries \eqref{eq:entry} contain the full logarithm $L$, whose coefficients of high order have large
denominators. The first lemma shows that a fixed truncation of $L$ leaves $\Delta$ unchanged.

Weight the monomial $T^sV^\beta$ by $v_0s+W(\beta)/\theta$. Let $\fa_H$ be the closed ideal of $\C[[T,V]]$
spanned by the monomials of weight at least $H$. The quotient $\C[[T,V]]/\fa_H$ has the basis $\cJ_H$.

\begin{lemma}[{\cite[Lemma~15.1]{ChudnovskyChudnovsky26}}]\label{lem:trunc}
Let $T_1,\dots,T_m$ be positive integers with $T_iv_0\ge w_i/\theta$, and let $L_i^{\rm tr}$ be the sum of the
terms of $L$ of degree less than $T_i$. Let $A^{\rm tr}$ be the matrix \eqref{eq:entry} with $L$ replaced by
$L_i^{\rm tr}$ in the factor with index $i$. Then $\det A^{\rm tr}=\det A$ for every choice of the columns.
\end{lemma}

\begin{proof}
Put $R_i=L-L_i^{\rm tr}$ and let $\Phi$ be the substitution $V_i\mapsto V_i+R_i(T)$. Every term of $R_i$ has
degree at least $T_i$, so its weight is at least $T_iv_0\ge w_i/\theta$, the weight of $V_i$. Hence $\Phi$ and
its inverse $V_i\mapsto V_i-R_i(T)$ do not decrease weight. They preserve $\fa_H$ and induce inverse linear maps
of $\C[[T,V]]/\fa_H$. The map $\Phi$ carries the truncated evaluation at each centre to the full one. Order the
basis $\cJ_H$ by $|\beta|$. The image of $T^sV^\beta$ is $T^sV^\beta$ plus terms of smaller $|\beta|$, so the
matrix of $\Phi$ is unitriangular. Thus each centre block of $A$ is a matrix of determinant $1$ times the
corresponding block of $A^{\rm tr}$.
\end{proof}

\begin{proposition}\label{prop:arith}
Let $F_0\ge1/\theta$ and suppose $\Delta\ne0$. Then
\[
 \frac{\log|\Delta|}{MH}\ge-(1-\bb)-\Ear,\qquad \Ear=\frac{c_0F_0}{v_0},\qquad c_0=1.03883.
\]
\end{proposition}

This is \cite[Lemma~15.2 and Corollary~15.3]{ChudnovskyChudnovsky26} with $\Z$ in place of the order
$\Z[i\sqrt d]$, and with the constant of Rosser and Schoenfeld in place of $\log4+4$. We give the proof.

\begin{proof}
Put $T_i=\lceil F_0w_i/v_0\rceil$ and $\ell_i=\lcm(1,\dots,T_i-1)$, with $\ell_i=1$ if $T_i\le2$. Then
$T_iv_0\ge F_0w_i\ge w_i/\theta$, so Lemma~\ref{lem:trunc} gives $\Delta=\det A^{\rm tr}$. By
Proposition~\ref{prop:chart}, $U\in\Z[[T]]$ and $\ell_iL_i^{\rm tr}\in\Z[T]$. Hence
\[
 q_i\ell_i\Bigl(j\frac{p_i}{q_i}+L_i^{\rm tr}\Bigr)=\ell_ijp_i+q_i\ell_iL_i^{\rm tr}\in\Z[T].
\]
By \eqref{eq:entry}, $\prod_i(q_i\ell_i)^{\alpha_{ci}-\beta_{ri}}A^{\rm tr}_{rc}$ is an integer whenever
$\beta_r\le\alpha_c$, and $A^{\rm tr}_{rc}=0$ otherwise.

Put $d_i=\sum_c\alpha_{ci}-\sum_r\beta_{ri}$. Since $\Delta\ne0$, some permutation $\varphi$ of the Leibniz
expansion gives a nonzero product. Then $\beta_r\le\alpha_{\varphi(r)}$ for all $r$, and
$d_i=\sum_r(\alpha_{\varphi(r)i}-\beta_{ri})\ge0$. The same identity holds for every permutation. So
$\Delta\prod_i(q_i\ell_i)^{d_i}$ is a sum of integers. It is a nonzero integer, and
\[
 \log|\Delta|\ge-\sum_id_i(\log q_i+\log\ell_i).
\]
Now $\log q_i\le w_i$. The number $\log\lcm(1,\dots,n)$ is Chebyshev's function at $n$, and Rosser and
Schoenfeld proved that it is less than $c_0n$ for $n\ge1$ \cite[Theorem~12]{RosserSchoenfeld62}. Hence
$\log\ell_i\le c_0(T_i-1)<c_0F_0w_i/v_0$. By \eqref{eq:S},
$\sum_iw_id_i=\sum_cW(\alpha_c)-\sum_rW(\beta_r)\le MH(1-\bb)$. Therefore
\[
 \log|\Delta|\ge-MH(1-\bb)\Bigl(1+\frac{c_0F_0}{v_0}\Bigr)\ge-MH\bigl((1-\bb)+\Ear\bigr).\qedhere
\]
\end{proof}

The error $\Ear$ depends on $F_0$ and $v_0$ only. It does not involve $K$ or the grid.

%% =====================================================================
\section{The analytic upper bound}\label{sec:analytic}

The upper bound follows \cite[Section~16]{ChudnovskyChudnovsky26}. The grid enters at two places: the transport
of the determinant uses $j\Omega\in\Lambda$, and the collision estimate uses a disc containing all centres. Put
\[
 D=\Bigl\lceil\frac H{2w_0}\Bigr\rceil,\qquad \delta_{ji}=c_{ji}-j\Omega=j\Bigl(\frac{p_i}{q_i}-\Omega\Bigr).
\]
By \eqref{eq:S}, $k_c+2b_c\le D$ for every column.

\subsection{Transport}

\begin{lemma}\label{lem:transport}
Let $\tA$ be the matrix with entries
\[
 \tA_{(j,s,\beta),c}=[z^s]\,\tg_{D,k_c,b_c}(j\Omega+z)\binom{\alpha_c}\beta\prod_i
 \bigl(j\Omega+z+\delta_{ji}\bigr)^{\alpha_{ci}-\beta_i}.
\]
Then $|\det\tA|=|\Delta|\,|\psi(0)|^{DJ_HK}$.
\end{lemma}

\begin{proof}
Fix a centre $j$ and fix $\beta$. The set of $s$ with $(s,\beta)\in\cJ_H$ is an initial segment of $\N$.
Substitute $T=\tau(z)$ in the series \eqref{eq:packet}. The coefficient of $z^sV^\beta$ is a combination of the
coefficients of $T^{s'}V^\beta$ with $s'\le s$, with diagonal coefficient $\tau'(0)^s=(-1)^s$. This is a row
operation of determinant $\pm1$ on the block of $j$.

By Proposition~\ref{prop:germ} the substituted series is
$u(z)^{k}t(z)^{b}\prod_i(c_{ji}+V_i+z)^{\alpha_i}$. Since $j\Omega\in\Lambda$, the germs of $u$ and $t$ at
$0$ and at $j\Omega$ agree. Also $c_{ji}+z=j\Omega+z+\delta_{ji}$. Taking the coefficient of $V^\beta$ gives
$u(j\Omega+z)^kt(j\Omega+z)^b\binom\alpha\beta\prod_i(j\Omega+z+\delta_{ji})^{\alpha_i-\beta_i}$.

Multiply this by $\tpsi(j\Omega+z)^D=\tpsi(z)^D$, which is a series with constant term $\psi(0)^D$. On the
coefficients of $z^s$ this is again lower triangular, with diagonal $\psi(0)^D$. Over the block of $j$ it
multiplies the determinant by $\psi(0)^{DJ_H}$. The product is
$\tg_{D,k,b}(j\Omega+z)\binom\alpha\beta\prod_i(j\Omega+z+\delta_{ji})^{\alpha_i-\beta_i}$. The $K$ blocks
give the identity.
\end{proof}

The corresponding identity of \cite[Lemma~16.1]{ChudnovskyChudnovsky26} has the factor
$\prod_j|\psi(\gamma_j)|^{DJ_H}$. For $\gamma_j=j\Omega$, Lemma~\ref{lem:mult}(a) gives $\psi(j\Omega)=e^{j^2\eta_\Omega\Omega}\psi(0)$. The periodised multiplier
removes this Gaussian in $j$.

\subsection{Expansion in the displacements}
Let $\cT_H=\{\xi\in\N^m:W(\xi)\le H\}$. For a column $c$ and $\xi\in\cT_H$ define the entire function
\[
 G_{c,\xi}(\zeta)=\binom{\alpha_c}\xi\,\tg_{D,k_c,b_c}(\zeta)\,\zeta^{|\alpha_c|-|\xi|},
\]
which is $0$ unless $\xi\le\alpha_c$. For a row with centre $j$ and transverse index $\beta$ put
$s(\beta,\xi,j)=\binom\xi\beta\prod_i\delta_{ji}^{\xi_i-\beta_i}$. The identity
\[
 \binom{\alpha_i}{\beta_i}(\zeta+\delta)^{\alpha_i-\beta_i}
 =\sum_{\beta_i\le\xi_i\le\alpha_i}\binom{\xi_i}{\beta_i}\delta^{\xi_i-\beta_i}\binom{\alpha_i}{\xi_i}
 \zeta^{\alpha_i-\xi_i}
\]
with $\zeta=j\Omega+z$ expresses each row of $\tA$ as a combination of rows of entire functions. Since
$W(\alpha_c)\le H$, only $\xi\in\cT_H$ occur. Multilinearity in the rows gives the finite identity
\begin{equation}\label{eq:tags}
 \det\tA=\sum_{g}\Bigl(\prod_rs(\beta_r,g(r),j(r))\Bigr)\det B_g,
\end{equation}
over all maps $g$ from the rows to $\cT_H$, where $(B_g)_{r,c}=[z^{s_r}]\,G_{c,g(r)}(j(r)\Omega+z)$ and
$r=(j(r),s_r,\beta_r)$.

\subsection{Collision}
The next lemma is the collision estimate of \CM. It uses only that the same entire functions are expanded at all
the centres, and that the centres lie in a disc of half the radius.

\begin{lemma}[{\cite[Lemma~16.2]{ChudnovskyChudnovsky26}}]\label{lem:collision}
Let $R>0$ with $(K-1)\Omega+1\le R/2$, and suppose $|G_{c,\xi}(\zeta)|\le e^{H\cG}$ for $|\zeta|\le R$ and all
$c,\xi$. For a map $g$ put $n_\xi=|g^{-1}(\xi)|$. Let $\kappa=(\log2)/4$, $s_0=2^{-1/2}$ and
$\rho_H=(\log M+\kappa-\log(1-s_0))/H$. Then
\[
 |\det B_g|\le\exp\Bigl(-\kappa\sum_\xi n_\xi^2+MH(\cG+\rho_H)\Bigr).
\]
\end{lemma}

\begin{proof}
Put $a_{c,\xi,l}=R^l[\zeta^l]G_{c,\xi}$, so $|a_{c,\xi,l}|\le e^{H\cG}$ by Cauchy's inequality. Put
$b_{r,l}=[z^{s_r}]((j(r)\Omega+z)/R)^l$. Cauchy's inequality on $|z|=1$ gives
$|b_{r,l}|\le((j(r)\Omega+1)/R)^l\le2^{-l}$. So
$(B_g)_{r,c}=\sum_{l\ge0}b_{r,l}a_{c,g(r),l}$, with absolute convergence. Expanding the determinant by
multilinearity in the rows,
\[
 \det B_g=\sum_{(l_r)}\Bigl(\prod_rb_{r,l_r}\Bigr)\det\bigl[a_{c,g(r),l_r}\bigr]_{r,c}.
\]
Two rows with equal pairs $(g(r),l_r)$ are equal, so the determinant vanishes. In a surviving term the $l_r$ are
distinct within each class $g^{-1}(\xi)$, and $\sum_rl_r\ge\sum_\xi\binom{n_\xi}2$. Each coefficient
determinant is at most $M!\,e^{MH\cG}$. Write $2^{-\sum l_r}=s_0^{\sum l_r}s_0^{\sum l_r}$. Bound the first
factor by $s_0^{\sum_\xi\binom{n_\xi}2}$ and sum the second over all $(l_r)$. This gives
\[
 |\det B_g|\le M!\,e^{MH\cG}s_0^{\sum_\xi\binom{n_\xi}2}(1-s_0)^{-M}.
\]
Finally $M!\le M^M$, $\sum_\xi n_\xi=M$ and $\log s_0=-2\kappa$.
\end{proof}

\subsection{The estimate}
Put
\[
 r_*=2(\Omega+1),\qquad C_\Lambda=\log S_*+\log^+\frac1{|\psi(0)|}+\chi(r_*+Y)^2.
\]
These constants depend only on $\Lambda$, $a$ and the choice of $z_Q$.

\begin{theorem}\label{thm:analytic}
Assume $\Delta\ne0$. Let $\nu\ge1$ with $|\Omega-p_i/q_i|\le q_i^{-\nu}$ for $1\le i\le m$, and let
$0<w_{\min}\le\min_iw_i$. Fix $a_0,\gamma\in(0,1)$ and put
\[
 N_H=|\{\xi\in\N^m:W(\xi)\le a_0H\}|,\qquad c_H=\frac{\kappa\gamma^2M}{HN_H}.
\]
Then
\[
 \frac{\log|\Delta|}{MH}\le\Ean+e_H+\max\bigl\{-c_H,\ -\nu\bigl(a_0(1-\gamma)-\bb\bigr)\bigr\},
\]
where
\[
 \Ean=\frac{C_\Lambda K^2}{2w_0}+\frac{\nu+\log(4r_*K^2)}{w_{\min}},\qquad
 e_H=\frac{C_\Lambda K^2}H+\rho_H+\frac{\log|\cT_H|}H.
\]
\end{theorem}

\begin{proof}
\emph{Displacements.} Since $\log q_i>w_i-1$, we have $|\delta_{ji}|\le Kq_i^{-\nu}\le\exp(C_\delta-\nu w_i)$
with $C_\delta=\nu+\log K>0$. A nonzero $s(\beta,\xi,j)$ has $\beta\le\xi$, and $|\xi|\le H/w_{\min}$ for
$\xi\in\cT_H$. Since $\binom\xi\beta\le2^{|\xi|}$,
\begin{equation}\label{eq:disp}
 |s(\beta,\xi,j)|\le\exp\Bigl(-\nu\bigl(W(\xi)-W(\beta)\bigr)+\frac{H(\log2+C_\delta)}{w_{\min}}\Bigr).
\end{equation}

\emph{The disc.} Take $R=r_*K$. Then $(K-1)\Omega+1\le(\Omega+1)K=R/2$. On $|\zeta|\le R$ we have
$|\operatorname{Im}\zeta|\le R$ and $R+Y\le(r_*+Y)K$. Lemma~\ref{lem:growth}, $\binom{\alpha_c}\xi\le2^{|\alpha_c|}$
and $|\alpha_c|\le H/w_{\min}$ give $|G_{c,\xi}|\le e^{H\cG}$ with
\[
 \cG=\frac DH\bigl(\log S_*+\chi(r_*+Y)^2K^2\bigr)+\frac{\log(2R)}{w_{\min}}.
\]

\emph{Two cases.} Consider a map $g$ whose scalar product in \eqref{eq:tags} is nonzero. Then
$W(g(r))\ge W(\beta_r)$ for every row, so $\sum_rW(g(r))\ge MH\bb$. Suppose first that
$\sum_rW(g(r))\ge a_0(1-\gamma)MH$. The product of the bounds \eqref{eq:disp} contains the factor
$\exp(-\nu MH(a_0(1-\gamma)-\bb))$, and we use Lemma~\ref{lem:collision} without its first term. Otherwise
more than $\gamma M$ rows have $W(g(r))\le a_0H$. Indeed, if at least $(1-\gamma)M$ rows had $W(g(r))>a_0H$,
the sum would be at least $a_0(1-\gamma)MH$. These rows use at most $N_H$ values of $g$, so the
Cauchy--Schwarz inequality gives $\sum_\xi n_\xi^2\ge\gamma^2M^2/N_H$. Lemma~\ref{lem:collision} then gives the
factor $\exp(-c_HMH)$, and the factor $\exp(-\nu\sum_r(W(g(r))-W(\beta_r)))$ is at most $1$. There are
$|\cT_H|^M$ maps $g$. Hence
\[
 \frac{\log|\det\tA|}{MH}\le\frac{\log2+C_\delta}{w_{\min}}+\cG+\rho_H+\frac{\log|\cT_H|}H
 +\max\bigl\{-c_H,-\nu(a_0(1-\gamma)-\bb)\bigr\}.
\]

\emph{The multiplier.} By Lemma~\ref{lem:transport} and $MH=KJ_HH$,
\[
 \frac{\log|\Delta|}{MH}\le\frac{\log|\det\tA|}{MH}+\frac DH\log^+\frac1{|\psi(0)|}.
\]
The terms with the factor $D/H$ add up to at most $(D/H)C_\Lambda K^2$, since $K\ge1$. Moreover
$D/H\le1/(2w_0)+1/H$. Finally $\log2+C_\delta+\log(2R)=\nu+\log(4r_*K^2)$.
\end{proof}

\begin{corollary}[{\cite[Corollary~16.4]{ChudnovskyChudnovsky26}}]\label{cor:limits}
With all data other than $H$ fixed, $e_H\to0$ and
\[
 c_H\to c_\infty=\frac{\kappa\gamma^2K\theta^m}{(m+1)v_0a_0^m}\qquad(H\to\infty).
\]
\end{corollary}

\begin{proof}
The number of points of $\N^n$ in the simplex $\{\sum_ia_ix_i<X\}$ is asymptotic to $X^n/(n!\prod_ia_i)$. Hence
$J_H\sim\theta^mH^{m+1}/((m+1)!\,v_0\prod_iw_i)$ and $N_H\sim a_0^mH^m/(m!\prod_iw_i)$. This gives the limit
of $c_H=\kappa\gamma^2KJ_H/(HN_H)$. The numbers $M$ and $|\cT_H|$ grow polynomially in $H$, so
$e_H\to0$.
\end{proof}

%% =====================================================================
\section{The parameters}\label{sec:params}

In this section $\nu>3$ is rational. We choose the remaining parameters explicitly. Put
\[
 b=\frac56-\frac1{2\nu},\qquad h=\frac{\nu-3}{4\nu},\qquad\theta=1-h,
\]
\[
 a_0=1-bh,\qquad\gamma=\frac{h(\nu-3)}{12\nu},\qquad F_0=\frac1\theta,
\]
and
\[
 g=\nu\bigl(a_0(1-\gamma)-\theta\bigr)-(1-\theta).
\]

\begin{lemma}\label{lem:abc}
\begin{enumerate}[label=(\alph*)]
\item $2/3<b<1-1/\nu$ and $0<h<1/4$. Hence $0<\theta<a_0<1$.
\item $a_0^3<\theta^2$.
\item $g\ge h(\nu-3)/12>0$.
\end{enumerate}
\end{lemma}

\begin{proof}
(a) We have $3b-2=(\nu-3)/(2\nu)=2h>0$ and $1-1/\nu-b=(\nu-3)/(6\nu)>0$. Since $b<1$, $a_0>\theta$.

(b) Expanding,
$a_0^3-\theta^2=h\bigl((2-3b)+(3b^2-1)h-b^3h^2\bigr)$. With $2-3b=-2h$ the bracket equals
$3(b^2-1)h-b^3h^2<0$.

(c) We have $a_0-\theta=h(1-b)$ and $\nu(1-b)-1=(\nu-3)/6$. Hence
$g=h(\nu-3)/6-\nu\gamma a_0\ge h(\nu-3)/6-\nu\gamma=h(\nu-3)/12$.
\end{proof}

By (b), $a_0^{3/2}<\theta$, that is, $a_0/\theta<a_0^{-1/2}$. Choose rationals $c_1,c_2$ with
\[
 \frac{a_0}\theta<c_1<a_0^{-1/2},\qquad a_0<c_2<\min\{c_1^{-2},\ c_1\theta\}.
\]
The second interval is not empty, because $c_1<a_0^{-1/2}$ means $a_0<c_1^{-2}$ and $c_1>a_0/\theta$ means
$a_0<c_1\theta$. Also $c_1>1$ and $c_1\theta<a_0^{-1/2}\theta<a_0^{1/2}<1$. For integers $m\ge1$ put
\begin{equation}\label{eq:choice}
 K=\lfloor c_1^m\rfloor,\qquad w_0=c_2^{-m},\qquad v_0=4K\theta^mw_0.
\end{equation}
These are rational. Let $\Ear$ be as in Proposition~\ref{prop:arith} and $c_\infty$ as in
Corollary~\ref{cor:limits}.

\begin{proposition}\label{prop:params}
As $m\to\infty$:
\begin{enumerate}[label=(P\arabic*)]
\item $K\theta^m\to0$;
\item $K(w_0/v_0)\theta^m=1/4$ for every $m$;
\item $K^2/w_0\to0$;
\item $v_0\to\infty$, so $\Ear\to0$;
\item $c_\infty\to\infty$.
\end{enumerate}
\end{proposition}

\begin{proof}
For $c_1^m\ge2$ we have $c_1^m/2\le K\le c_1^m$. (P1) $K\theta^m\le(c_1\theta)^m$. (P2) is the definition of
$v_0$. (P3) $K^2/w_0\le(c_1^2c_2)^m$ and $c_1^2c_2<1$. (P4) $v_0\ge2(c_1\theta/c_2)^m$ and $c_1\theta>c_2$.
(P5) Substituting $v_0$, the factor $K$ cancels and
$c_\infty=\kappa\gamma^2(c_2/a_0)^m/(4(m+1))$, with $c_2>a_0$.
\end{proof}

The factor $K$ cancels from $c_\infty$ because $v_0$ is proportional to $K$. This is why the analytic cost
$K^2/w_0$ and the collision rate $c_\infty$ meet in the single condition $a_0^3<\theta^2$.

%% =====================================================================
\section{Proof of Theorems~\ref{cor:curve} and~\ref{thm:main}}\label{sec:proof}

\begin{theorem}\label{thm:omega}
For every $\nu>3$ there are only finitely many pairs $(p,q)$ of integers with $q\ge1$ and
$|\Omega-p/q|\le q^{-\nu}$.
\end{theorem}

\begin{proof}
Decreasing $\nu$ only enlarges the set of pairs, so we may assume $\nu$ rational. Suppose that infinitely many
pairs satisfy the inequality. Their denominators are unbounded. Since $\Omega\ne0$, all but finitely many have
$p\ne0$ and $q\ge2$. Take the parameters of Section~\ref{sec:params} and the constant $C_\Lambda$ of
Theorem~\ref{thm:analytic}.

\emph{Step 1.} By Proposition~\ref{prop:params} we may fix $m$ so large that \eqref{eq:epa-base} holds,
$c_\infty>2$, and
\[
 \Ear+\frac{C_\Lambda K^2}{2w_0}<\min\Bigl(\frac g4,\frac14\Bigr).
\]
This fixes $K$, $w_0$, $v_0$ and $\theta$. Choose $\sigma_0$, $C_4$ and $q_*$ as in Theorem~\ref{thm:epa}.

\emph{Step 2.} Choose pairs $(p_1,q_1),\dots,(p_m,q_m)$ from the given infinite set in turn. Each $w_i$ is to
satisfy \eqref{eq:epa-thresholds}, which involves only $w_1,\dots,w_{i-1}$. Each is also to satisfy $w_i\ge w_*$,
where
\[
 \frac{\nu+\log(4r_*K^2)}{w_*}<\min\Bigl(\frac g4,\frac14\Bigr).
\]
Unbounded denominators allow every choice. Put $w_{\min}=w_*$. Then $E=\Ear+\Ean<\min(g/2,1/2)$.

\emph{Step 3.} Theorem~\ref{thm:epa} gives $D_{\mathrm{EPA}}$. The centres $c_{ji}=jp_i/q_i$ are injective in
$j$ for each $i$, so there is $H_0$ such that for every $H\ge H_0$ divisible by $D_{\mathrm{EPA}}$ some choice
of columns gives $\Delta_H\ne0$. Let $\bb_H\in[0,\theta)$ be the corresponding mean.

\emph{Step 4.} Proposition~\ref{prop:arith} and Theorem~\ref{thm:analytic} apply to the same $\Delta_H$:
\begin{align*}
 \frac{\log|\Delta_H|}{MH}&\ge-(1-\bb_H)-\Ear,\\
 \frac{\log|\Delta_H|}{MH}&\le\Ean+e_H+\max\bigl\{-c_H,-\nu(a_0(1-\gamma)-\bb_H)\bigr\}.
\end{align*}
If the maximum is $-c_H$, subtracting the lower bound from the upper bound gives
$0\le E+e_H-c_H+1-\bb_H\le1+E+e_H-c_H$. By Corollary~\ref{cor:limits}, $e_H\to0$ and $c_H\to c_\infty>2$,
while $E<1/2$. So the right side is negative for large $H$. In the other case the difference is
\[
 E+e_H-\bigl(\nu a_0(1-\gamma)-1-(\nu-1)\bb_H\bigr)\le E+e_H-g,
\]
because $\nu>1$ and $\bb_H<\theta$. Since $E<g/2$, this is negative for large $H$. Both cases are impossible
for large $H$ divisible by $D_{\mathrm{EPA}}$.
\end{proof}

\begin{proof}[Proof of Theorem~\ref{thm:main}]
By Lemma~\ref{lem:integral} and the remarks after it, $\omega=N\Omega$ with $N$ a nonzero integer, for the
lattice $\Lambda=d^{-1}L$. Let $\varepsilon>0$. If infinitely many pairs satisfied
$|\omega-p/q|\le q^{-3-\varepsilon}$, Lemma~\ref{lem:scale} with $\nu'=3+\varepsilon/2$ would contradict
Theorem~\ref{thm:omega}. So the finitely many exceptions have bounded $q$, and $q_0$ exists. If $\omega$ were
a rational $a/b$, the pairs $(na,nb)$ would give infinitely many exceptions. So $\omega$ is irrational, and
Dirichlet's theorem gives $\mu(\omega)\ge2$.
\end{proof}

\begin{proof}[Proof of Theorem~\ref{cor:curve}]
Take a Weierstrass equation $y^2+a_1xy+a_3y=x^3+a_2x^2+a_4x+a_6$ of $E$ over $\Q$. With
$y_1=y+(a_1x+a_3)/2$ and $x_1=x+(a_2+a_1^2/4)/3$ it becomes $y_1^2=x_1^3+A_1x_1+B_1$ with $A_1,B_1\in\Q$, and
$dx/(2y+a_1x+a_3)=dx_1/(2y_1)$. The invariant differentials defined over $\Q$ are the rational multiples of this
differential, so $\eta=r\,dx_1/(2y_1)$ with $r\in\Q^\times$. By \cite[Section~VI.5]{Silverman09} the
periods of $dx_1/(2y_1)$ form the lattice $\Lambda_1$ with $g_2(\Lambda_1)=-4A_1$ and $g_3(\Lambda_1)=-4B_1$. The
periods of $\eta$ form $r\Lambda_1$, whose invariants $r^{-4}g_2(\Lambda_1)$ and $r^{-6}g_3(\Lambda_1)$ are
rational. Theorem~\ref{thm:main} applies to $L=r\Lambda_1$.
\end{proof}

%% =====================================================================
\section{Why the exponent is three}\label{sec:three}

The exponent $3$ comes from the analytic term $C_\Lambda K^2/(2w_0)$. This section explains how that term
determines the exponent, why it arises, and what would remove it.

\subsection{The threshold for a cost of order \texorpdfstring{$K^e$}{K\textasciicircum e}}
Suppose that the analytic estimate of Theorem~\ref{thm:analytic} held with $CK^e/w_0$ in place of
$C_\Lambda K^2/(2w_0)$, for some $e\ge1$. Take $K=\lfloor c_1^m\rfloor$, $w_0=c_2^{-m}$ and
$v_0=CK\theta^mw_0$ with fixed $c_1,c_2,C$, as in \eqref{eq:choice}. Up to the constant this choice of $v_0$ is the best
one. The interpolation theorem needs $v_0>K\theta^mw_0$, and $c_\infty$ decreases as $v_0$ grows. The proof
of Section~\ref{sec:proof} needs four rates. The analytic term must tend to $0$, which needs $c_1^ec_2\le1$. The
collision rate $c_\infty$, proportional to $(c_2/a_0)^m/(m+1)$, must tend to infinity, which needs $c_2>a_0$.
The arithmetic error $c_0F_0/v_0$ must stay small, which needs $v_0$ bounded below and so $c_1\theta\ge c_2$. Finally
$g>0$ needs $\nu(a_0-\theta)>1-\theta$. The first three give $a_0<c_2\le\min(c_1^{-e},c_1\theta)$, hence
$a_0/\theta<c_1<a_0^{-1/e}$ and
\begin{equation}\label{eq:window}
 a_0^{e+1}<\theta^e.
\end{equation}

\begin{proposition}\label{prop:threshold}
Let $e\ge1$ be real and $\nu>1$. There are $0<\theta<a_0<1$ with $a_0^{e+1}<\theta^e$ and
$\nu(a_0-\theta)>1-\theta$ if and only if $\nu>e+1$.
\end{proposition}

\begin{proof}
Write $\theta=x^{e+1}$ with $0<x<1$. Condition \eqref{eq:window} says $a_0<x^e$. Since $1-\theta>0$, the
second condition requires $a_0>\theta$ and $\nu>(1-\theta)/(a_0-\theta)$. For $\theta<a_0<x^e$,
\[
 \frac{1-\theta}{a_0-\theta}>r_e(x),\qquad r_e(x)=\frac{1-x^{e+1}}{x^e-x^{e+1}}.
\]
The function $k_e(x)=1-(e+1)x^e+ex^{e+1}$ vanishes at $x=1$, and $k_e'(x)=e(e+1)x^{e-1}(x-1)<0$ on
$(0,1)$. So $k_e>0$ on $(0,1)$, which says $r_e(x)>e+1$. Hence both conditions force $\nu>e+1$.
Conversely, let $\nu>e+1$. By l'H\^opital's rule $r_e(x)\to e+1$ as $x\to1$, so we can fix $x$ with
$r_e(x)<\nu$. By continuity some $a_0$ slightly below $x^e$ satisfies both conditions.
\end{proof}

When \eqref{eq:window} holds, choose $c_1\in(a_0/\theta,a_0^{-1/e})$, $c_2\in(a_0,\min(c_1^{-e},c_1\theta))$
and $\gamma$ small. Then all four rates hold strictly, as in Proposition~\ref{prop:params}. So within this family
of parameters the exponent obtained is exactly $e+1$. For $e=1$ the window is $a_0^2<\theta$. This is the
condition of \cite[Theorem~17.1]{ChudnovskyChudnovsky26} and of the proof for $\pi$ \cite{OpenAIPi2026}, and it
gives the exponent $2$. For $e=2$ it is Lemma~\ref{lem:abc}(b), and it gives $3$.

\subsection{The growth of the multiplier along the grid}
The cost $K^2$ is a property of the disc. The collision estimate expands the same entire functions at all the
centres, so it needs a disc of radius $R$ of order $K\Omega$ around the origin. On that disc the multiplier
attains the growth $\exp(D\chi R^2)$ near $\pm iR$, by Lemma~\ref{lem:mult}(b). No other multiplier does
better by more than a constant factor in the exponent.

\begin{remark}\label{rem:jensen}
Let $\Phi$ be entire with $\Phi(0)\ne0$, vanishing to order at least $D$ at every point of $\pm z_Q+\Lambda$.
Such a function is needed to clear the poles of $u^D$. Let $n(r)$ count these points in $|z|<r$. Jensen's formula
gives
\[
 \log\max_{|z|=R}|\Phi|\ge\log|\Phi(0)|+D\int_0^R\frac{n(r)}r\,dr.
\]
Comparing areas, $n(r)\ge2\pi r^2/(\Omega Y)-O(r+1)$. Hence
$\log\max_{|z|=R}|\Phi|\ge\log|\Phi(0)|+D\pi R^2/(\Omega Y)-O(D(R+1))$. With $R$ of order $K$ and $D/H$ close to $1/(2w_0)$,
any Cauchy estimate on the disc costs a term of order $K^2/w_0$.
\end{remark}

The entries of $\tA$ themselves are not large. By periodicity $\tg_{D,k,b}(j\Omega+z)=\tg_{D,k,b}(z)$, so
the entries are bounded on the strip $|\operatorname{Im}z|\le1$ independently of $j$, up to the polynomial
factors. The loss comes from bounding the Taylor coefficients at $0$ through the maximum on $|\zeta|=R$.

\subsection{A collision estimate without that growth}
Suppose the analytic estimate held with a cost $CK/w_0$, for a constant $C$ depending only on $\Lambda$. By
Proposition~\ref{prop:threshold} with $e=1$, the proof of Section~\ref{sec:proof} would then give
$\mu(\omega)=2$ for every nonzero real $\omega$ in every lattice with rational invariants. Such an estimate
would have to separate the rows of different centres without expanding at a single point. The monomials
$\zeta^l$ achieve this separation in Lemma~\ref{lem:collision}: rows of the same tag at different centres share
a coefficient vector only when they share the exponent $l$. A basis adapted to the strip, on which translation by
$\Omega$ acts diagonally, would keep the entries bounded. It would need a new counting argument for coincident
rows. We do not know such an estimate.

\subsection{Algebraic centres and norms}
A two-dimensional grid of periods with algebraic multipliers is available only under complex multiplication,
by Remark~\ref{rem:noCM}. Other algebraic centres bring a different loss. Suppose the cleared determinant is a
nonzero algebraic integer of a number field $k$ of degree $n$. The lower bound then comes from
$|N_{k/\Q}(\Delta)|\ge1$, and it needs upper bounds for all conjugates of $\Delta$. The conjugates are
determinants at the conjugate centres, which are not close to periods. They obey only the trivial bound, without
the collision or displacement gains. The arithmetic side then loses a factor of order $n$. Under complex
multiplication with the rectangular grid, the only other embedding of $\Q(i\sqrt d)$ is complex conjugation,
and $|a+bi\sqrt d|^2$ is a positive integer, so nothing is lost. For the powers $\varpi^n$ of the lemniscate
period the natural centres involve $(p/q)^{1/n}$, of degree $n$. There G.~and D.~Chudnovsky use the norm from
$\Q(i,(p/q)^{1/n})$ to $\Q(i)$, and they obtain $\mu(\varpi^n)=O(n\log2n)$ \cite[Theorem~18.1]{ChudnovskyChudnovsky26}.

%% =====================================================================
\section*{Disclosure of AI use}
The mathematics and the text of this paper were produced by AI systems at the request of the author. The author
posed the problem, directed the work and is responsible for the paper. The systems were Claude Opus~5.5, Claude
Sonnet~5.5 and Claude Fable~5.2 (Anthropic), running in Claude Code. The construction, the nonvanishing theorem
and the method of the proof are due to G.~V.~Chudnovsky and D.~V.~Chudnovsky \cite{Chudnovsky26,ChudnovskyChudnovsky26},
building on OpenAI's proof for $\pi$ \cite{OpenAIPi2026}.

\providecommand{\bysame}{\leavevmode\hbox to3em{\hrulefill}\thinspace}
\providecommand{\MR}{\relax\ifhmode\unskip\space\fi MR }
% \MRhref is called by the amsart/book/proc definition of \MR.
\providecommand{\MRhref}[2]{%
  \href{http://www.ams.org/mathscinet-getitem?mr=#1}{#2}
}
\providecommand{\href}[2]{#2}
\begin{thebibliography}{10}

\bibitem{Chudnovsky26}
Gregory~V. Chudnovsky, \emph{The irrationality exponent of the lemniscate
  period: a determinant proof with a {L}ean 4 formalization}, 2026, Preprint,
  19 October 2026, \href{https://arxiv.org/abs/2610.26539}{arXiv:2610.26539}.

\bibitem{ChudnovskyChudnovsky26}
Gregory~V. Chudnovsky and David~V. Chudnovsky, \emph{Irrationality exponents of
  real {CM} elliptic periods and powers of the lemniscate period}, 2026,
  Preprint, 22 October 2026,
  \href{https://arxiv.org/abs/2610.31112}{arXiv:2610.31112}.

\bibitem{David95}
Sinnou David, \emph{Minorations de formes lin\'eaires de logarithmes
  elliptiques}, M\'em. Soc. Math. France (N.S.) \textbf{62} (1995), 1--143.

\bibitem{Feldman51}
Naum~I. Fel'dman, \emph{The approximation of certain transcendental numbers.
  {II}. {T}he approximation of certain numbers connected with the {W}eierstrass
  function {$\wp(z)$}}, Izv. Akad. Nauk SSSR Ser. Mat. \textbf{15} (1951),
  no.~2, 153--176, In Russian. MR 42471, Zbl 0044.27105.

\bibitem{Feldman74}
\bysame, \emph{The periods of elliptic functions}, Acta Arith. \textbf{24}
  (1974), no.~5, 477--489, In Russian; English title as in Waldschmidt08.

\bibitem{HirataKohno90}
Noriko Hirata-Kohno, \emph{Mesures de transcendance pour les quotients de
  p\'eriodes d'int\'egrales elliptiques}, Acta Arith. \textbf{56} (1990),
  no.~2, 111--133.

\bibitem{HirataKohno91}
\bysame, \emph{Formes lin\'eaires de logarithmes de points alg\'ebriques sur
  les groupes alg\'ebriques}, Invent. Math. \textbf{104} (1991), no.~2,
  401--433.

\bibitem{LiuJiangZhang26}
Jingwen Liu, Kai Jiang, and Pingwen Zhang, \emph{Irrationality exponents of
  logarithms of positive rational numbers}, 2026, Preprint, 7 October 2026,
  revised 5 November 2026 (v3),
  \href{https://arxiv.org/abs/2610.10192}{arXiv:2610.10192}.

\bibitem{Masser75}
David~W. Masser, \emph{Elliptic functions and transcendence}, Lecture Notes in
  Math., vol. 437, Springer-Verlag, Berlin, 1975.

\bibitem{DLMF}
{NIST}, \emph{Digital {L}ibrary of {M}athematical {F}unctions}, Chapter 23,
  \url{https://dlmf.nist.gov/23.2}, 2026.

\bibitem{OpenAIPi2026}
{OpenAI}, \emph{The irrationality exponent of {$\pi$} is 2}, Preprint dated
  September 24, 2026, released October 6, 2026, with a Lean formalization,
  \href{https://github.com/openai/math/tree/fd4aeeb2ee4fc729c18d98444fed42fd0529eeeb/preprints/The-irrationality-exponent-of-pi-is-2-September-24-2026}{paper}
  and
  \href{https://github.com/openai/math/tree/fd4aeeb2ee4fc729c18d98444fed42fd0529eeeb/lean/OAI/NumberTheory/PiExponent}{Lean
  formalization} (repository commit \texttt{fd4aeeb}), 2026.

\bibitem{PhilipponWaldschmidt88}
Patrice Philippon and Michel Waldschmidt, \emph{Formes lin\'eaires de
  logarithmes sur les groupes alg\'ebriques commutatifs}, Illinois J. Math.
  \textbf{32} (1988), no.~2, 281--314.

\bibitem{PhilipponWaldschmidt89}
\bysame, \emph{Formes lin\'eaires de logarithmes elliptiques et mesures de
  transcendance}, Th\'eorie des nombres ({Q}u\'ebec, 1987), de Gruyter, Berlin,
  1989, pp.~798--805.

\bibitem{RosserSchoenfeld62}
J.~Barkley Rosser and Lowell Schoenfeld, \emph{Approximate formulas for some
  functions of prime numbers}, Illinois J. Math. \textbf{6} (1962), no.~1,
  64--94.

\bibitem{Roth55}
Klaus~F. Roth, \emph{Rational approximations to algebraic numbers}, Mathematika
  \textbf{2} (1955), 1--20.

\bibitem{Schneider37}
Theodor Schneider, \emph{Arithmetische {U}ntersuchungen elliptischer
  {I}ntegrale}, Math. Ann. \textbf{113} (1937), no.~1, 1--13.

\bibitem{Siegel32}
Carl~Ludwig Siegel, \emph{{\"U}ber die {P}erioden elliptischer {F}unktionen},
  J. Reine Angew. Math. \textbf{167} (1932), 62--69.

\bibitem{Silverman09}
Joseph~H. Silverman, \emph{The arithmetic of elliptic curves}, second ed.,
  Grad. Texts in Math., vol. 106, Springer, Dordrecht, 2009.

\bibitem{Waldschmidt08}
Michel Waldschmidt, \emph{Elliptic functions and transcendence}, Surveys in
  Number Theory, Dev. Math., vol.~17, Springer, New York, 2008, pp.~143--188.

\end{thebibliography}

\end{document}
